\documentclass[10pt,a4paper]{article}
\usepackage[margin=0.35in]{geometry}
\usepackage{fontspec}
\usepackage{xcolor}
\usepackage[hidelinks]{hyperref}

\setmainfont{Noto Sans CJK SC}
\XeTeXlinebreaklocale "zh"
\XeTeXlinebreakskip = 0pt plus 1pt
\setlength{\parindent}{0pt}
\setlength{\parskip}{0pt}
\emergencystretch=3em
\sloppy
\pagestyle{empty}

\definecolor{sectiongray}{HTML}{333333}

\newcommand{\cvsection}[1]{%
  \vspace{0.4em}{\large\bfseries\color{sectiongray}#1}\par
  \vspace{-0.45em}\rule{\linewidth}{0.4pt}\vspace{0.1em}
}

\let\olditemize\itemize
\let\endolditemize\enditemize
\renewenvironment{itemize}{%
  \olditemize
  \setlength{\itemsep}{0pt}
  \setlength{\topsep}{0pt}
  \setlength{\parsep}{0pt}
  \setlength{\partopsep}{0pt}
  \setlength{\leftmargin}{1.25em}
}{\endolditemize}

\newcommand{\entry}[4]{%
  \textbf{#1} \hfill {\small #2}\\
  \if\relax\detokenize{#3}\relax
  \else
    {\small #3}%
    \if\relax\detokenize{#4}\relax\else \hfill {\small #4}\fi
    \par
  \fi
  \vspace{-0.2em}
}

\begin{document}
\fontsize{9}{10.1}\selectfont

\begin{center}
{\Large \textbf{王中原}}\\[-0.2em]
{\small \href{mailto:wang1504626@gmail.com}{wang1504626@gmail.com} \quad | \quad \href{https://www.linkedin.com/in/zhongyuanwang}{linkedin.com/in/zhongyuanwang}}\\[0.2em]
{\small 创始机器学习科学家，专注 AI、LLM、Agent 和无监督图模型。20 项专利，5 家企业客户覆盖美国、印度和海湾地区。}
\end{center}
\vspace{-0.5em}

\cvsection{工作经历}
\entry{创始机器学习科学家}{2024 至今}{RaptorX.AI，Remote}{}
\begin{itemize}
\item 构建 criteria-driven multi-agent autopilots，覆盖论文与专利生产: 文献/gap 扫描、实验核验、TMLR gates、prior-art search、专利 drafting、counsel handoff 和 audit log。
\item 从零构建无监督欺诈检测引擎，结合自研 GNN、行为 autoencoder、可解释性和 LLM 调查报告，统一无监督检测、图欺诈团伙发现与 audit-ready explainability。
\item 设计数千个特征，覆盖 4 个业务场景，处理 1000 万+ 交易、59 万+ 实体和 280 万条边。基于 30 类关系的 edge-probability model 降低对欺诈标签的依赖。
\item 实现 328$\times$ 候选压缩 (590K $\rightarrow$ 1,883 high-confidence alerts)，欺诈类型 specificity 达 99.9\%。
\item 将 PoC 转化为 5 家签约企业客户: Payactiv (US)、NSDL 与 Lenskart (India)、Omantel 与 Thawani Pay (Oman)。美国部署覆盖 46 万+ 用户和 6,700+ 雇主。
\end{itemize}

\entry{高级工程师，自动驾驶 AI 与数据平台}{2021 至 2024}{Lotus Tech (Geely Group)，杭州}{}
\begin{itemize}
\item 负责车端采集数据到数据中心的全自研视频匿名化 pipeline，支撑 80+ TB/day 数据流。YOLOv5 模型在人脸与车牌上达到 90\%+ recall，覆盖 1000 万+ 帧，满足 GDPR/PIPL 合规。
\item 通过 GPU 加速 inference 将延迟降低 30\%，在车队规模替代第三方供应商，节省约 \$1M/year。
\end{itemize}

\entry{独立机器学习学习与 Graph-ML 项目}{2020 至 2021}{}{}
\begin{itemize}
\item 系统学习 graph machine learning，在公开 Kaggle 数据集上实现 GNN node classification 与 link prediction。
\end{itemize}

\entry{软件开发工程师，金融反欺诈}{2020}{DataVisor，上海}{}
\begin{itemize}
\item 为头部商业银行构建无监督贷款申请欺诈模型，结合 GNN 与传统 ML，使用行为、设备、第三方征信和原始交易数据的数百个特征。
\item 通过技术主导和客户演示，直接贡献 \$500,000 企业合同落地。
\end{itemize}

\entry{研究实习生，UCL-DiDi Master Research Program}{2019}{DiDi AI Labs，北京}{}
\begin{itemize}
\item 项目: Leveraging Graph Neural Networks for User Travel Intent Prediction。导师: Prof. Jieping Ye (DiDi Chuxing VP, IEEE Fellow)。
\item 作为 UCL cohort 30+ 候选项目中仅 2 个入选 DiDi AI Labs research program 的项目之一，为后续在 RaptorX 将 GNN 应用于金融欺诈检测奠定基础。
\end{itemize}

\cvsection{专利与论文}
\textbf{专利}: 发明人，20 项申请，覆盖图欺诈检测与 LLM/agent systems。当前状态: 6 项美国已 filed，4 项印度已 published，10 项 draft。印度 published app nos.: 202541024623, 202541045354, 202441079562, 202541032797。

\textbf{论文}: 第一作者，均 under review。
\begin{itemize}
\item When the Tool Decides: LLM Agents Defer Blindly to GNN Tools (TMLR; \href{https://arxiv.org/abs/2606.14476}{arXiv:2606.14476})
\item LLM Features Can Hurt GNNs: Concatenation Interference on Homophilous Graph Benchmarks (TMLR; \href{https://arxiv.org/abs/2606.17579}{arXiv:2606.17579})
\item How Much Do Deep GNNs Actually Help? Decomposing Depth, Features, and Structure (NeurIPS 2026 E\&D)
\end{itemize}

\cvsection{技术技能}
\textbf{ML / Graph}: PyTorch, PyTorch Geometric (PyG), GNNs, autoencoders, anomaly detection, link prediction, fraud-ring detection, LLM-based investigation\\
\textbf{Data \& MLOps}: Spark / PySpark, Kafka, ClickHouse, Redis, Parquet, Docker, Kubernetes, ArgoCD, CI/CD, FastAPI / gRPC, MLflow\\
\textbf{Cloud}: AWS (S3 / EC2 / ECR), Grafana, Prometheus, Loki

\cvsection{教育经历}
\entry{University College London, UK}{2018 至 2019}{MSc in Data Science and Machine Learning}{}
Thesis: Graph Convolutional Networks for user behavior prediction (jointly with DiDi AI Labs)。NLP 小组研究项目成绩位列 300 名学生前 3\%。\\[0.2em]
\entry{King{'}s College London, UK}{2015 至 2018}{BSc in Computer Science (Artificial Intelligence)}{}

\end{document}
